% =============================================================================
% SECTION 5 — ANALYSIS (revised for paper_v2)
% Headline table kept inline; the patch-shape figure is moved to the appendix.
% Trajectory-level statistics are run-means over three independent evaluation
% runs of each checkpoint on SWE-Bench-Verified (500 instances per run).
% =============================================================================

\section{Analysis}
\label{sec:analysis}
\vspace{-0.5ex}

The previous section established that FIM mid-training improves end-task
metrics. We now ask \emph{how} the resulting agent behaves differently
and \emph{where} along the trajectory the gains accrue. Both analyses
use the 14B configuration with R2E-Gym, comparing the
post-training-only baseline (\emph{R2E-Gym}) against our recipe
(\emph{FIM-Midtrain $+$ R2E-Gym}); statistics are run-means over three
evaluation runs on SWE-Bench-Verified. The Lite analogue, the full
action-type distribution, the no-patch mechanism, and a concrete
trajectory contrast are deferred to Appendix~\ref{app:behavior_extra}.

\vspace{-0.5ex}
\subsection{Recovery from Negative Observations}
\label{sec:analysis:recovery}

A trajectory contains a \emph{negative observation} if any tool output
matches a fixed set of error patterns (stack traces,
``\texttt{No replacement was performed}'', shell errors, etc.; full list
in Appendix~\ref{app:algo}). The fraction of such trajectories is
essentially identical across checkpoints ($88.8\%$ baseline vs.\
$91.8\%$ ours), so the agents see comparable amounts of negative
feedback. The \emph{recovery rate}---the fraction of error-containing
trajectories that nonetheless terminate with a passing patch---rises
from $24.8\%$ to $28.8\%$ ($+4.0$~pp; $+16\%$ relative;
Table~\ref{tab:behavior_main}). This is precisely the capability our
framing predicts mid-training should support: continuing productively
after an external return contradicts the model's prior expectation.
Mid-training also shifts the agent toward an iterate-and-verify policy,
increasing edits per solved task from $3.3$ to $7.4$ and trajectory
length from $15.1$ to $23.6$ steps; action-type breakdowns are in
Appendix~\ref{app:behavior_extra}.

% -----------------------------------------------------------------------------
% Compact headline table
% -----------------------------------------------------------------------------
\begin{table}[h]
\centering
\caption{Headline trajectory metrics on SWE-Bench-Verified (14B, R2E-Gym),
run-means over three runs. Full metrics in Appendix~\ref{app:behavior_extra}.}
\label{tab:behavior_main}
\small
\setlength{\tabcolsep}{8pt}
\vspace{-0.5ex}
\begin{tabular}{lcccc}
\toprule
Setting & Recovery rate (\%) & Edits / solved & Steps / solved & Pass (\%) \\
\midrule
+ R2E-Gym                          & 24.8           & 3.3           & 15.1           & 26.2 \\
\textbf{+ FIM-Midtrain + R2E-Gym}  & \textbf{28.8}  & \textbf{7.4}  & \textbf{23.6}  & \textbf{29.2} \\
\bottomrule
\end{tabular}
\vspace{-1ex}
\end{table}

\vspace{-0.5ex}
\subsection{The Gain Concentrates on Multi-Function Reasoning}
\label{sec:analysis:multifn}

The most direct test of the isomorphism argument
(Section~\ref{sec:method:motivation}) is whether mid-training
preferentially helps tasks that require reasoning about cross-function
dependencies inside a file. We stratify the $500$ Verified tasks by the
shape of the gold patch (Figure~\ref{fig:patchshape},
Appendix~\ref{app:behavior_extra}). On the $88$ tasks whose gold patch
modifies $\geq 2$ functions within a single file, the baseline solves
$13.6\%$ on average and ours solves $22.7\%$, an absolute gain of
$+9.1$~pp---more than $4{\times}$ the gain on the $341$ single-function
tasks ($+2.1$~pp). Per-instance head-to-head on this bucket shows ours
uniquely solves about twice as many tasks as the baseline. Multi-file
tasks ($n{=}71$) are not differentially helped, which we attribute to a
granularity mismatch: our FIM operates within files
(Appendix~\ref{app:behavior_extra}). The slice where mid-training helps
most is precisely the slice where the agent must follow control- and
data-dependencies between caller-callee or sibling functions inside a
file---the same structure exposed by function-aware FIM masking. A
complementary outcome-distribution breakdown
(Figure~\ref{fig:failure_app}) shows the $+15$-task gain is driven by an
order-of-magnitude reduction in \emph{no-patch} failures alongside small
reductions in localization errors; the FIM signal conditions the model
to produce a non-empty span between prefix and suffix---a disposition
that survives subsequent post-training.
